\section*{अध्याय \ref{ch:b2:vertebrate-development} --- कशेरुकियों का भ्रूणीय विकास}

\begin{solution}{exo:b2:vertebrate-development:1}
\omterm{def:b2:vertebrate-development:egg}{विदलन}: बिना वृद्धि के तेज़ समसूत्री विभाजन, जो \omterm{def:b2:vertebrate-development:blastula}{कोरकपुटी} बनाते हैं।
\omterm{def:b2:vertebrate-development:blastula}{कोरकपुटी}: \omterm{def:b2:vertebrate-development:blastula}{कोरकगुहा} के चारों ओर छोटी कोशिकाओं की खोखली गेंद या
चक्रिका। गैस्ट्रुलन: वे कोशिका-गतियाँ जो मध्यचर्म तथा
अंतश्चर्म को भीतर ले आती हैं और आँत-गुहा तथा पृष्ठरज्जु सहित
तीन-परती गैस्ट्रुला बनाती हैं। तंत्रिकायन: पृष्ठीय
बाह्यचर्म का मुड़कर \omterm{prop:b2:vertebrate-development:neurulation}{तंत्रिका नली} बनना, कायखंडों के खंडीभवन सहित,
जो कशेरुकी शरीर-योजना वाला न्यूरुला बनाता है।
\end{solution}

\begin{solution}{exo:b2:vertebrate-development:2}
बाह्यचर्म: अधिचर्म, मस्तिष्क तथा मेरुरज्जु, परिधीय तंत्रिकाएँ
(\omterm{prop:b2:vertebrate-development:neurulation}{तंत्रिका शिखा}), लेंस और वर्णक कोशिकाएँ। मध्यचर्म: पृष्ठरज्जु, पेशी,
अस्थि तथा उपास्थि, वृक्क, हृदय तथा रक्त, चर्म। अंतश्चर्म:
आँत का अस्तर, यकृत, अग्न्याशय, फेफड़े, अवटु ग्रंथि।
\end{solution}

\begin{solution}{exo:b2:vertebrate-development:3}
कम \omterm{def:b2:vertebrate-development:egg}{पीतक}: पूरा, लगभग बराबर \omterm{def:b2:vertebrate-development:egg}{विदलन} और किसी खोखली गेंद में
गोल कोरकरंध्र से अंतर्वलन द्वारा गैस्ट्रुलन (मेंढक:
पूरा पर असमान, जिसमें पीतकी वनस्पति कोशिकाएँ बड़ी और धीमी होती हैं)।
अधिक \omterm{def:b2:vertebrate-development:egg}{पीतक} (चूज़ा): \omterm{def:b2:vertebrate-development:egg}{विदलन} \omterm{def:b2:vertebrate-development:egg}{पीतक} के ऊपर एक चक्रिका तक सीमित, और
गैस्ट्रुलन उस चपटी चक्रिका में एक दरार, यानी आद्य रेखा, से होकर,
और परतें किसी गुहा को घेरने के बजाय \omterm{def:b2:vertebrate-development:egg}{पीतक} पर फैलती हुई।
\end{solution}

\begin{solution}{exo:b2:vertebrate-development:4}
\omterm{def:b2:vertebrate-development:induction}{प्रेरण} किसी सक्षम ऊतक की नियति का किसी पड़ोसी ऊतक के संकेत से
बदल जाना है। वनस्पति कोशिकाएँ अपने ऊपर के सीमांत क्षेत्र में
मध्यचर्म प्रेरित करती हैं; \omterm{def:b2:vertebrate-development:induction}{संघटक} (पृष्ठीय ओष्ठ का मध्यचर्म) पृष्ठीय
बाह्यचर्म में तंत्रिका पट्ट प्रेरित करता है; और नेत्र पुटिका
सिर के बाह्यचर्म में लेंस प्रेरित करती है।
\end{solution}

\begin{solution}{exo:b2:vertebrate-development:5}
$2^{12} = 4096$ कोशिकाएँ, जिनकी त्रिज्या $1/2^{4} = \qty{62.5}{\micro m}$;
और पृष्ठ $2^{4} = 16$ गुना। वह अतिरिक्त पृष्ठ उन उपकलाओं की
झिल्ली है जिन्हें गैस्ट्रुलन मोड़कर त्वचा तथा आँत बनाएगा, और वह
क्षेत्रफल भी जिससे \omterm{def:b2:vertebrate-development:blastula}{कोरकपुटी} अपने परिवेश से विनिमय करती है।
\end{solution}

\begin{solution}{exo:b2:vertebrate-development:6}
$10^{-4}\times 2^{k} = 0.4$: $2^{k} = 4000$, $k = 12$। चार गुना
DNA के साथ वह अनुपात $4\times 10^{-4}$ से आरंभ होता है और $0.4$ तक
$2^{k} = 1000$ पर पहुँच जाता है: यानी दसवाँ विभाजन, दो पहले।
\end{solution}

\begin{solution}{exo:b2:vertebrate-development:7}
$k = \ln 2/600 = \qty{1.16e-3}{s^{-1}}$; $\lambda = \sqrt{2\times
10^{-12}/1.16\times 10^{-3}} = \qty{42}{\micro m}$। देहलियाँ
$x = \lambda\ln 2 = \qty{29}{\micro m}$ और $\lambda\ln 10 =
\qty{96}{\micro m}$ पर: यानी 29, 67 और शेष की पट्टियाँ।
\end{solution}

\begin{solution}{exo:b2:vertebrate-development:8}
स्रोत दुगुना करने पर दोनों सीमाएँ $\lambda\ln
2 = \qty{29}{\micro m}$ बाहर खिसक जाती हैं: पहली पट्टी दुगुनी हो जाती है
और दूसरी अपनी चौड़ाई बनाए रखती है। $k$ दुगुना करने पर $\lambda$ $\sqrt 2$
से भाग जाता है (\qty{29}{\micro m}) और, चूँकि $c_0 = j/\sqrt{Dk}$ है,
$c_0$ भी $\sqrt 2$ से भाग जाता है: सीमाएँ $\lambda'\ln(c_0'/T) =
29\times(0.69 - 0.35) = \qty{10}{\micro m}$ और $29\times(2.30 -
0.35) = \qty{57}{\micro m}$ पर चली जाती हैं --- यानी दोनों पट्टियाँ सिकुड़ती हैं।
\end{solution}

\begin{solution}{exo:b2:vertebrate-development:9}
किसी अरंजित न्यूट गैस्ट्रुला का पृष्ठीय ओष्ठ किसी रंजित के
अधर भाग पर ग्राफ़्ट किया गया; वह भीतर लुढ़का और पोषी ने एक
दूसरा भ्रूण बनाया --- \omterm{prop:b2:vertebrate-development:neurulation}{तंत्रिका नली}, कायखंड, आँत --- जो पहले से जुड़ा था।
रंजकता के अंतर ने दिखाया कि वह दूसरी अक्ष पोषी की कोशिकाओं की
बनी थी, इसलिए उस ग्राफ़्ट ने उन्हें प्रेरित किया था, स्वयं वह अक्ष नहीं
बनाई थी। इसने स्पष्टीकरण के रूप में ग्राफ़्ट के स्व-विभेदन को ख़ारिज
कर दिया, और दिखाया कि कोई छोटा क्षेत्र वह संकेत ढोता है जो पूरी अक्ष
को संघटित करता है।
\end{solution}

\begin{solution}{exo:b2:vertebrate-development:10}
(क) पेट की त्वचा: आरंभिक बाह्यचर्म सक्षम तो है पर अभी प्रेरित नहीं,
इसलिए वह अपने नए परिवेश का अनुसरण करता है। (ख) तंत्रिका ऊतक: पश्च
गैस्ट्रुला तक वह प्रेरित हो चुका है और निर्धारित है। (ग) तंत्रिका पट्ट:
पश्च-गैस्ट्रुला के पेट का बाह्यचर्म अब भी सक्षम है, और पृष्ठरज्जु उसे
प्रेरित कर देता है। (घ) कोई पृष्ठीय संरचना नहीं: पेट के ऊतक की एक गेंद,
क्योंकि न कोई मध्यचर्म को पृष्ठीय नियतियों की ओर प्रेरित करता है, न
बाह्यचर्म को तंत्रिका नियतियों की ओर।
\end{solution}

\begin{solution}{exo:b2:vertebrate-development:11}
\qty{18}{\celsius} पर: $90 + 11\times 30 = \qty{420}{min}$, यानी 7 घंटे।
\qty{10}{\celsius} पर: \qty{17.5}{h}। वह संक्रमण दोनों स्थितियों में
बारहवें विभाजन पर पड़ता है, क्योंकि उसका चालक DNA और कोशिकाद्रव्य का
अनुपात है, जो विभाजनों की संख्या पर निर्भर है और इस पर नहीं कि उनमें
कितना समय लगा: वह घड़ी दुगुनेपन गिनती है, घंटे नहीं।
\end{solution}

\begin{solution}{exo:b2:vertebrate-development:12}
नियामक विकास में कोशिकाएँ अपनी नियतियाँ पड़ोसियों के \omterm{def:b2:vertebrate-development:induction}{प्रेरण} से पाती
हैं, इसलिए वह भ्रूण हानि के बाद फिर गढ़ सकता है (बँटे अंडाणु से
जुड़वाँपन, ग्राफ़्ट किए क्षेत्र का पुनर्जनन); जबकि मोज़ेक
विकास में नियतियाँ स्थानीकृत कोशिकाद्रव्यी निर्धारकों के साथ विरासत में
मिलती हैं, इसलिए हर कोरकखंड केवल अपना भाग बनाती है और बँटा
भ्रूण दो आधे लार्वा बनाता है। वह भेद मात्रा भर का है: मेंढक का
धूसर अर्धचंद्र एक निर्धारक है, और मोज़ेक भ्रूण बाद में \omterm{def:b2:vertebrate-development:induction}{प्रेरण} काम में
लेते हैं; हर भ्रूण विरासत और बातचीत मिलाता है, और कशेरुकी बातचीत पर
अधिक देर टिके रहते हैं।
\end{solution}

\begin{solution}{pb:b2:vertebrate-development:1}
\textbf{1.} अंडाणु $\tfrac43\pi(0.7)^{3} = \qty{1.44}{mm^3}$; केंद्रक
$\tfrac43\pi(0.03)^{3} = \qty{1.1e-4}{mm^3}$; अनुपात $8\times
10^{-5}$।
\textbf{2.} $2^{k} = 4000$: $k = 12$, यानी 4096 कोशिकाएँ।
\textbf{3.} $75 + 11\times 30 = \qty{405}{min}$, यानी लगभग 6.75 घंटे।
\textbf{4.} त्रिज्या $0.7/2^{4} = \qty{44}{\micro m}$; अनुपात $8\times
10^{-5}\times 4096 = 0.32$ --- यानी किसी साधारण कोशिका का: केंद्रक
कोशिकाद्रव्य के बराबर आ चुके हैं।
\textbf{5.} $2^{4} = 16$।
\textbf{6.} लगभग $4095\,c$ DNA; और उस अंडाणु में बारह विभाजनों लायक़
हिस्टोन, पॉलिमरेज़ तथा न्यूक्लियोटाइड सँजोए हैं, इसलिए कोई
अनुलेखन नहीं चाहिए।
\textbf{7.} युग्मनजी जीन चालू हो जाते हैं, विभाजन धीमे पड़ते हैं और
समकालिकता खो देते हैं, कोशिकाएँ गतिशील हो जाती हैं। अतिरिक्त DNA इंजेक्ट
करने पर वह संक्रमण उतने ही विभाजन पहले आ जाता है जितना DNA अधिक है,
जिसे विभाजनों या काल की कोई गिनती समझा ही नहीं सकती।
\textbf{8.} शुक्राणु-प्रवेश बिंदु के सामने, उस धूसर अर्धचंद्र पर जो
निषेचन के बाद अंडाणु-बाह्यक के घूमने से तय हुआ।
\textbf{9.} पहले पृष्ठीय मध्यचर्म (पूर्वरज्जुकी पट्ट, फिर पृष्ठरज्जु)
मध्यरेखा में आद्यांत्र की छत पर; उसके बग़ल में पार्श्व मध्यचर्म
(कायखंड, पार्श्व पट्ट); अंतश्चर्म
आद्यांत्र को अस्तरित करता है; और बाह्यचर्म अधिवर्धन से बाहर की ओर फैल जाता है।
\textbf{10.} $\qty{2000}{\micro m}/\qty{600}{min} = \qty{3.3}{\micro
m/min}$: यानी किसी तंतुकोरक का तिगुना --- कोई समन्वित चादर अकेली
कोशिका से तेज़ चलती है।
\textbf{11.} उन कोशिका-गतियों को युग्मनजी जीन-उत्पाद (नए
आसंजन अणु, गतिशीलता) और धीमे चक्र चाहिए; और उस संक्रमण से पहले वे
कोशिकाएँ केवल विभाजित ही हो सकती हैं।
\textbf{12.} हर रज्जुकी के पास किसी न किसी अवस्था में पृष्ठरज्जु होता
है, और कशेरुकी उसी के चारों ओर कशेरुक-दंड गढ़ते हैं; \omterm{prop:b2:vertebrate-development:neurulation}{तंत्रिका नली}
उसी से प्रेरित होती है और उस योजना का परिणाम है, उद्गम नहीं।
\textbf{13.} बोतल-कोशिकाओं के संकुचन के बिना कोई कोरकरंध्र नहीं बनता,
कुछ भीतर नहीं लुढ़कता, और वह भ्रूण अपनी परतें अपनी जगह रखे एक गेंद ही बना
रहता है: न आँत, न पृष्ठरज्जु, न कोई अक्ष।
\textbf{14.} $k = \ln 2/1200 = \qty{5.8e-4}{s^{-1}}$; $\lambda =
\sqrt{10^{-12}/5.8\times 10^{-4}} = \qty{42}{\micro m}$।
\textbf{15.} $x_1 = 42\ln 2.5 = \qty{38}{\micro m}$; $x_2 = 42\ln 12.5
= \qty{105}{\micro m}$।
\textbf{16.} \qty{38}{\micro m} (2.5 कोशिकाएँ), \qty{67}{\micro
m} (4.5 कोशिकाएँ) और शेष की पट्टियाँ।
\textbf{17.} $x_1 = 42\ln 1.25 = \qty{9}{\micro m}$; $x_2 = 42\ln 6.25
= \qty{77}{\micro m}$: दोनों सीमाएँ $\lambda\ln 2 =
\qty{29}{\micro m}$ भीतर खिसक जाती हैं; पहली पट्टी 38 से घटकर 9 रह जाती
है, बीच वाली अपनी चौड़ाई बनाए रखती है, और बाहरी बढ़ जाती है।
\textbf{18.} पढ़े जाने से पहले ले जाई जाए तो वह नियति A बनती है ---
वह उसी सांद्रता को पढ़ती है जहाँ वह अब है; और बाद में ले जाई जाए तो वह
नियति B रखती है --- वह निर्धारित हो चुकी थी।
\textbf{19.} $x_T = \lambda\ln(c_0/T)$: स्रोत में $f$ गुणक का
परिवर्तन हर सीमा को $\lambda\ln f$ खिसकाता है, इसलिए दो गुने की कोई
त्रुटि उस प्रतिरूप को केवल $0.7\lambda$ खिसकाती है, जो कई $\lambda$
फैले किसी क्षेत्र का एक अंश भर है --- यानी वह प्रतिरूप आकारजन की ठीक-ठीक
मात्रा के प्रति लगभग असंवेदी है।
\textbf{20.} पृष्ठीय ओष्ठ अधर भाग पर: एक दूसरी अक्ष। अधर सीमांत
क्षेत्र पृष्ठीय भाग पर: वह अपने परिवेश से फिर से निर्दिष्ट होकर
पृष्ठीय ऊतक बनाता है; और कुछ अतिरिक्त नहीं।
\textbf{21.} किसी न्यूरुला में वह बाह्यचर्म अब सक्षम नहीं रहा: वह
ग्राफ़्ट पृष्ठरज्जु तो बनाता है पर कोई दूसरी \omterm{prop:b2:vertebrate-development:neurulation}{तंत्रिका नली} प्रेरित नहीं करता।
\textbf{22.} अर्धचंद्र के तल के साथ बँटने पर: दो पूरे भ्रूण;
और लंब बँटने पर: एक भ्रूण और पेट के ऊतक की एक गेंद।
\textbf{23.} अकेले: अधिचर्मी गेंद। वनस्पति कोशिकाओं के साथ: उस टोपी में
प्रेरित मध्यचर्म --- पेशी, पृष्ठरज्जु।
\textbf{24.} यह सिद्ध करने के लिए कि वह दूसरी अक्ष पोषी कोशिकाओं में
प्रेरित हुई थी और ग्राफ़्ट ने गढ़ी नहीं थी; यदि वह ग्राफ़्ट अभेद्य होता
तो स्व-विभेदन ख़ारिज नहीं किया जा सकता था। किसी भिन्न जाति का ग्राफ़्ट
चलता है (उन्होंने न्यूट की दो जातियाँ काम में लीं), और यही
कोशिकाओं को पहचानने योग्य बनाता था।
\textbf{25.} बारहवाँ \omterm{def:b2:vertebrate-development:egg}{विदलन}, 4096 कोशिकाएँ, लगभग 6.75 घंटे पर; और
38 तथा \qty{67}{\micro m} की पट्टियाँ, और \qty{105}{\micro
m} से आगे सब कुछ।
\end{solution}
